\section{Quantitative approach and reporting}
We first describe each outcome on its own valid-response base. We then use separate models to describe adjusted associations in the recruited sample.
\subsection{Adjusted associations in the recruited sample}
The two core logit models (R0 and D0) relate reported remittance disposition and domestic digital participation, respectively, to IAT3, IADT2, IAFF2, age, grouped education and grouped recorded residence. Completed secondary education and Bogotá are the reference categories. The residual residence category combines Other, Soacha and Meta where those categories are documented; it is neither a fourth study city nor a sampling cluster. Two extensions (R1 and D1) add grouped occupation, with employees/workers as the reference, and the recorded arrival-year field (q2\_1). The question asks for first arrival, but for repeated entries its instruction specifies the most recent arrival to stay continuously. We do not use it to calculate tenure.

The core models use their own valid-outcome and complete-covariate samples. We check that the fitted and intended rows match. In the audited estimates, R0 and R1 use the same 402 records and D0 and D1 use the same 394 records, so each extension changes the controls without changing those rows. A fifth model, R0C, fits the remittance core on the intersection of the two core samples. That intersection equals the D0 sample, so no second domestic fit is needed. R0 versus R0C checks sensitivity to sample composition under the same remittance specification; it does not test whether the associations differ between outcomes.

For model $m$, let $S_m$ denote its estimation sample, and let $\widehat p_{im}(a)$ be the predicted probability for record $i$ when the focal measure is set to $a$ and the other covariates keep their observed values. We report the model-standardized difference
\[
  \widehat{\Delta}_m(a_0,a_1)=
  \frac{100}{|S_m|}\sum_{i\in S_m}
    [\widehat p_{im}(a_1)-\widehat p_{im}(a_0)] .
\]
Differences are in percentage points. We compare 0.50 with 0.75 for IAT3 and IADT2, and 0.50 with 1.00 for IAFF2. The IAFF2 comparison pools phone-only and proof-only at 0.50. The logit uses one slope for the score across 0, 0.50 and 1.00 and predicts the same probability for the two one-condition configurations at fixed covariates. It does not estimate associations for each configuration separately. The different score increments do not rank the measures. The technical supplement checks whether both exact endpoints occur within each education-by-recorded-residence stratum and whether either falls outside its observed range. This screen neither verifies overlap on all covariates nor changes the sample over which we average predictions.

The 95\% intervals are pointwise delta-method intervals using the model's independent-record sandwich covariance. They do not account for unrecorded referral or household dependence. We apply Holm's adjustment to one fixed family of six one-degree-of-freedom Wald tests: one coefficient-null test per focal measure in each core model \citep{holm1979}. The adjustment does not apply to the probability-contrast intervals. Current use may help respondents acquire the abilities they report. Neither covariate adjustment nor the common-sample comparison removes unmeasured confounding or selection into recruitment.

\subsection{Observed profiles}
We describe combinations of two reported onboarding conditions: own-name proof of address described as acceptable to financial institutions and an own-name line normally used for digital accounts or wallets. We count a respondent only when both answers are available and scorable. These combinations are observed answers, not fitted groups. The separate exploratory latent-class exercise below asks what a model gains by grouping responses; it does not assign classes in the main analysis.

\subsection{Survey results}
\label{sec:survey-results}
Domestic digital use was common: 368 of 413 respondents with valid answers (89.1\%) reported sending or receiving money digitally within Colombia in the previous 60 days. On the separate disposition question, 169 of 423 respondents (40.0\%) in the household-screened sample reported retaining or using some incoming proceeds digitally. The other valid answers reported all cash, not necessarily a cash withdrawal from a digital balance. The two rates are neither shares of money nor successive stages of a transfer.

\paragraph{Observed onboarding profiles.} Among the 404 respondents with scorable answers to both questions, 141 (34.9\%) reported neither condition, 128 (31.7\%) reported an own-name phone line only, 35 (8.7\%) reported own-name proof of address only, and 100 (24.8\%) reported both (Figure~\ref{fig:onboarding-conditions}). The other 19 respondents were unscored and are not counted as having neither condition. The regression contrast pools the two one-condition groups, although an own-name line alone was more than three times as common as own-name proof of address alone. These profiles do not establish which condition came first, whether a provider accepted a document, or whether a respondent controlled an account. The joint display (Table~\ref{tab:joint-evidence}) explains why these distinctions matter for interpreting the survey.

\begin{figure}[tbp]
\centering
\includegraphics[width=0.94\linewidth]{"\CFIDescriptivesDir F2_onboarding_conditions_artwork.pdf"}
\begin{minipage}{0.94\linewidth}
\caption{Reported own-name onboarding conditions}
\label{fig:onboarding-conditions}
\par\smallskip{\footnotesize\raggedright Notes: Unweighted shares among the 404 records with both answers scorable (19 records unscored). Respondents described whether their own-name address proof was acceptable to financial institutions. Phone refers to the number normally used for digital accounts or wallets. These answers do not establish verified eligibility, account ownership or control.\par Source: Authors' calculations from the CFI-DPI survey in Colombia.\par}
\end{minipage}
\end{figure}

\paragraph{Reported decision-making, information and access.} Among the 407 women with scorable answers on all five measurement items, 360 agreed that they usually decide how household remittances are used. Of those 360 women, 67.8\% (244) reported very little or no information about fees, exchange rates, and what to do if something goes wrong. On the same 407-case base, 63.9\% (200 of 313) of regular personal-smartphone users did not report own-name proof of address. Both comparisons use observed answers, not fitted profiles. The first does not establish who actually decides how to spend the money, and the second does not verify a provider's decision.

\paragraph{Adjusted associations.} Figure~\ref{fig:determinants-core} shows two positive core contrasts. In the remittance model, the average predicted probability is 16.35 percentage points higher at the two-condition score than at the pooled one-condition score (pointwise 95\% interval: 9.54 to 23.15). In the domestic model it is 8.15 points higher at 0.75 rather than 0.50 on the ability and confidence score (1.75 to 14.55). Both comparisons have exact endpoints in 11 of 12 education-by-residence strata, and one endpoint lies outside the observed range in one stratum for each comparison. The predictions retain all fitted respondents. The fraction counts strata; it gives neither the share of respondents affected nor overlap on all covariates.

The extension and common-sample estimates are reported with the core estimates in technical supplement Tables~2--3. The fixed family of six coefficient-null tests has two Holm-adjusted p-values below 0.05: the onboarding score in the remittance model and the ability and confidence score in the domestic model (technical supplement Table~5). This adjustment does not apply to the plotted intervals. Neither the estimates nor the tests compare effects across outcomes or tell us what changing documentation or training would do.

\begin{figure}[tbp]
\centering
\includegraphics[width=0.98\linewidth]{"\CFIDeterminantsDir det_contrasts_core.pdf"}
\caption{Adjusted associations with two digital-payment outcomes}
\label{fig:determinants-core}
\begin{minipage}{0.94\linewidth}
\smallskip\footnotesize\raggedright
Notes: R0 models reported digital retention or use of incoming remittances in the household-screened sample ($N=402$); D0 models reported domestic use ($N=394$). Points are higher-minus-lower predicted probability differences averaged over the full fitted sample, with pointwise 95\% intervals. IAT3 is the phone, internet and data score; IADT2, reported transfer ability and fraud-avoidance confidence; IAFF2, reported own-name phone registration and address proof. The endpoints are 0.50 and 0.75 for IAT3 and IADT2, and the pooled one-condition score and two-condition score for IAFF2. Open markers flag at least one stratum without both exact endpoints; a filled marker does not verify full overlap. The intervals do not account for unrecorded household or referral dependence.\par
Source: Authors' calculations from the CFI-DPI survey in Colombia.
\end{minipage}
\end{figure}

\subsection{What the exploratory latent profiles add}
\label{sec:exploratory-profiles-main}

The observed configurations and item tabulations remain the primary evidence. A separate exploratory six-item latent-class model (T02) describes four response profiles on 362 common cases with scorable answers across the candidate panels. These cases come from routes requiring regular phone use and answers on a reported qualifying account or wallet; the fit does not represent nonusers who did not reach those questions. Figure~\ref{fig:t02-response-profiles} shows fitted probabilities and mean memberships, not observed frequencies, assigned counts or population shares.

The profiles do not line up on one scale of inclusion. Profiles 1 and 4 have similar fitted probabilities of reporting high transfer ability (97.9\% and 98.9\%). Profile 1 is less likely to report an own-name phone line (29.8\% versus 82.3\%), but more likely to report regularly holding a balance for over a week (77.3\% versus 22.9\%). In Profile 4, holding a balance sometimes is the most likely answer (64.1\%).

Profile 2 is more likely than Profile 3 to report high transfer ability (69.1\% versus 6.7\%), but less likely to say one could set aside part of a remittance for saving or a business without permission (42.5\% versus 94.6\%). Profile 2 has the lowest fitted probabilities of regular balance holding (9.4\%) and always authenticating (4.5\%); Profile 3 has the highest probability of needing registration help (84.1\%). These are separate fitted responses, not observed combinations. The allocation question is hypothetical, not evidence of who actually decides.

With smoothing, four classes predict the six defining answers better than three on the internal folds. They do not improve prediction of the separate outcome of reported problems and trust endorsement: mean held-out log score falls from $-1.0625$ to $-1.0689$ and Brier loss rises from 0.5728 to 0.5766 (Table~\ref{tab:exploratory-profiles-main}, Panel B). Using the six coded responses directly, or as two multiple correspondence analysis scores, predicts that outcome better on both metrics under the same outcome model with residence. These are rankings on the same internal folds, not independent validation or tests of score differences. The technical supplement's \emph{Exploratory six-item profiles and class-free comparisons} section documents the recodes, smoothing and stability checks. We do not assign anyone to a fitted class.

\input{"\CFIPaperDir exhibits/exploratory-profiles.tex"}
\input{"\CFIPaperDir exhibits/t02-response-profiles.tex"}
\clearpage
